%=====================================================================
% T1 - METODOLOGIA CIENTIFICA - CEFET-RJ
% Modalidade 1: Planejamento de Estudo Primario
%
% ARQUIVO UNICO - pronto para colar no Overleaf.
% As referencias estao embutidas no bloco filecontents abaixo; ao
% compilar, o Overleaf gera o arquivo referencias.bib sozinho.
%
% COMO COMPILAR NO OVERLEAF
%   1. Novo projeto em branco, cole este arquivo inteiro no main.tex.
%   2. Menu > Compiler: pdfLaTeX.
%   3. Compile uma vez, depois "Recompile from scratch" (as citacoes so
%      aparecem apos a passagem do BibTeX).
%
% MARCADORES NO TEXTO
%   [PREENCHER]  -> dado, numero ou decisao sua
%   [VERIFICAR]  -> afirmacao que precisa de confirmacao na literatura
%
% A entrega final deve usar o template oficial da disciplina
% (main-planejamento-primario.tex). Transplante as secoes para la.
%=====================================================================

%---------------------------------------------------------------------
% REFERENCIAS EMBUTIDAS
%
% AVISO: nenhuma entrada abaixo e uma referencia real. Todas sao
% MARCADORES que descrevem o TIPO de fonte necessaria em cada ponto do
% texto. Substitua cada uma por um artigo real, localizado por voce,
% indexado no Scopus, e cujo conteudo voce tenha lido o suficiente para
% sustentar a afirmacao que o cita.
%
% O enunciado exige no minimo DEZ artigos indexados no Scopus.
% Ha 16 marcadores abaixo, para dar margem.
%
% FLUXO: busque no Scopus com as strings indicadas, exporte a entrada
% BibTeX direto da base, e substitua o marcador MANTENDO A CHAVE.
%---------------------------------------------------------------------

\begin{filecontents*}[overwrite]{referencias.bib}

@misc{ref:jogos_serios_reabilitacao,
  title = {[PREENCHER] Jogos serios aplicados a reabilitacao motora},
  note  = {Busca: ("serious game" OR gamification) AND rehabilitation
           AND (adherence OR engagement OR motivation)}
}

@misc{ref:biofeedback_engajamento,
  title = {[PREENCHER] Biofeedback e engajamento em reabilitacao},
  note  = {Se nao encontrar evidencia solida, ENFRAQUECA a afirmacao no
           texto em vez de forcar a citacao.
           Busca: biofeedback AND rehabilitation AND (engagement OR
           motivation OR adherence)}
}

@misc{ref:semg_fundamentos,
  title = {[PREENCHER] Fundamentos de eletromiografia de superficie},
  note  = {Pode ser artigo classico da area ou documento de
           padronizacao (ex.: recomendacoes SENIAM).
           Busca: "surface electromyography" AND (recommendations OR
           standards OR fundamentals)}
}

@misc{ref:semg_facial_clinico,
  title = {[PREENCHER] sEMG facial em contexto clinico},
  note  = {Idealmente um artigo que justifique o zigomatico maior como
           sitio de interesse.
           Busca: ("facial electromyography" OR "facial sEMG") AND
           (rehabilitation OR "facial palsy" OR asymmetry)}
}

@misc{ref:variabilidade_intersessao,
  title = {[PREENCHER] Variabilidade inter-sessao em sEMG},
  note  = {REFERENCIA MAIS IMPORTANTE: estabelece que o problema
           existe. Verifique ESPECIFICAMENTE se a literatura descreve a
           variabilidade como mudanca de ESCALA de amplitude -- essa e
           a premissa que sustenta o uso da distancia z-normalizada.
           Busca: ("surface electromyography" OR sEMG) AND
           ("inter-session" OR "session-to-session") AND (variability
           OR calibration OR recalibration)}
}

@misc{ref:fadiga_semg,
  title = {[PREENCHER] Fadiga muscular e alteracao do sinal de sEMG},
  note  = {ATENCAO: muitos artigos tratam de deslocamento ESPECTRAL,
           que exige sinal bruto. Como sua aquisicao entrega envelope,
           prefira fonte que trate de alteracao de AMPLITUDE.
           Busca: (sEMG OR electromyography) AND fatigue AND (amplitude
           OR "signal decline")}
}

@misc{ref:recuperacao_motora_semg,
  title = {[PREENCHER] Alteracao do padrao de sEMG ao longo da
           recuperacao motora},
  note  = {REFERENCIA-CHAVE PARA A SECAO 2.5. Sustenta que o sinal muda
           legitimamente com o progresso terapeutico -- e por isso o
           sistema precisa de plasticidade, e nao apenas de memoria.
           Busca: (sEMG OR electromyography) AND (recovery OR
           "motor learning" OR rehabilitation) AND (longitudinal OR
           "over time" OR progression)}
}

@misc{ref:matrix_profile_base,
  title = {[PREENCHER] Matrix Profile - formulacao original (STAMP)},
  note  = {Busca: "matrix profile" AND "time series" AND (motif OR
           discord)}
}

@misc{ref:matrix_profile_stomp,
  title = {[PREENCHER] STOMP - atualizacao incremental},
  note  = {Sustenta a escolha do STOMP pela viabilidade em operacao
           continua. Busca: STOMP AND "matrix profile" AND (incremental
           OR streaming)}
}

@misc{ref:matrix_profile_biosinais,
  title = {[PREENCHER] Matrix Profile aplicado a biosinais},
  note  = {Se existir aplicacao a EMG, e citacao obrigatoria: define se
           sua contribuicao e "aplicacao nova" ou "metodo novo". A
           ausencia tambem e um achado -- registre no texto.
           Busca: ("matrix profile" OR "motif discovery") AND
           (biosignal OR EMG OR physiological OR ECG)}
}

@misc{ref:aprendizado_continuo,
  title = {[PREENCHER] Aprendizado continuo - formulacao do problema},
  note  = {Base conceitual da Secao 2.5: dilema plasticidade x
           estabilidade, aprendizado sequencial sem acesso completo aos
           dados anteriores.
           Busca: ("continual learning" OR "lifelong learning") AND
           (stability OR plasticity OR sequential)}
}

@misc{ref:esquecimento_catastrofico,
  title = {[PREENCHER] Esquecimento catastrofico},
  note  = {Sustenta o problema que a consolidacao por replay resolve.
           Busca: "catastrophic forgetting" AND (neural OR incremental
           OR "continual learning")}
}

@misc{ref:replay_memoria,
  title = {[PREENCHER] Replay / memoria episodica em aprendizado
           continuo},
  note  = {Sustenta a Secao 4.4. Procure trabalhos sobre selecao de
           exemplares para o buffer de memoria (a critica ao sorteio
           uniforme e o que motiva a selecao por motifs).
           Busca: ("experience replay" OR "episodic memory" OR
           rehearsal) AND "continual learning" AND (exemplar OR
           selection)}
}

@misc{ref:deteccao_drift,
  title = {[PREENCHER] Deteccao de mudanca de regime / concept drift},
  note  = {Sustenta a Secao 4.5. Busca: ("concept drift" OR "change
           point") AND detection AND (stream OR "time series")}
}

@misc{ref:limiar_adaptativo,
  title = {[PREENCHER] Limiar adaptativo e deteccao de onset em sEMG},
  note  = {Busca: sEMG AND ("adaptive threshold" OR "onset detection")}
}

@misc{ref:jogo_semg_reabilitacao,
  title = {[PREENCHER] Jogo de reabilitacao controlado por sEMG},
  note  = {Trabalho relacionado direto. Analise explicitamente COMO ele
           trata calibracao: e dessa analise que sai sua lacuna.
           Busca: (sEMG OR electromyography) AND ("serious game" OR
           "video game") AND rehabilitation AND biofeedback}
}

@misc{ref:usabilidade_instrumento,
  title = {[PREENCHER] Instrumento de avaliacao de usabilidade},
  note  = {Verificar existencia de versao validada em portugues.
           Busca: (usability OR "user experience") AND (scale OR
           questionnaire) AND validation}
}

\end{filecontents*}

\documentclass[12pt,a4paper]{article}

\usepackage[utf8]{inputenc}
\usepackage[T1]{fontenc}
\usepackage[brazil]{babel}
\usepackage{amsmath,amssymb}
\usepackage{graphicx}
\usepackage{booktabs}
\usepackage{url}
\usepackage{indentfirst}
\usepackage[hidelinks]{hyperref}
\usepackage{tikz}
\usepackage{pgfgantt}
\usetikzlibrary{shapes.geometric, arrows.meta, positioning, calc}

\title{\textbf{EMGInvaders Adaptativo: Aprendizado Contínuo de Perfis de
Ativação Muscular para Calibração Persistente em um Jogo de
Reabilitação Baseado em Eletromiografia de Superfície Facial}}

\author{Matheus Kubrusly \and Renato Alves de Sousa \and Gustavo Guedes}
% [PREENCHER] Confirmar a composicao do grupo do T1. Limite de tres
% membros, e nao precisa coincidir com a autoria do EMGInvaders.

\date{Centro Federal de Educação Tecnológica Celso Suckow da Fonseca
(CEFET-RJ)}

\begin{document}
\maketitle

%---------------------------------------------------------------------
\begin{abstract}
\noindent
Jogos sérios controlados por eletromiografia de superfície (sEMG) são
usados para tornar sessões de reabilitação mais motivadoras, mas
esbarram em um obstáculo prático: o sinal muscular registrado muda de
uma sessão para outra, obrigando a repetir a calibração a cada uso.
Guardar simplesmente a calibração anterior, porém, não resolve o
problema, porque em reabilitação o sinal \textit{deve} mudar --- a
recuperação do paciente altera legitimamente seu padrão de ativação.
O sistema precisa, portanto, adaptar-se sem descartar o que já
aprendeu, e precisa distinguir a mudança que é ruído de medição
daquela que é resultado terapêutico. Este trabalho propõe tratar essa
questão como um problema de aprendizado contínuo. Propõe-se
representar cada participante por um conjunto de padrões recorrentes
de contração extraídos por Matrix Profile, atualizá-lo a cada sessão
com um mecanismo de consolidação que preserva o histórico, e ajustar o
limiar de ativação de forma contínua durante a sessão. A avaliação
planejada compara a proposta com a versão atual do EMGInvaders quanto
ao tempo de preparação, à retenção do perfil ao longo das sessões, à
estabilidade da detecção e à usabilidade percebida.

\vspace{0.5em}
\noindent\textbf{Palavras-chave:} eletromiografia de superfície;
reabilitação; aprendizado contínuo; Matrix Profile; jogos sérios.
\end{abstract}

%=====================================================================
\section{Introdução}
%=====================================================================

A recuperação de funções motoras após lesões neurológicas ou
periféricas depende da repetição frequente de exercícios específicos.
Essa repetição é justamente o que torna o processo cansativo, e a
desistência do paciente é um obstáculo conhecido ao sucesso do
tratamento~\cite{ref:jogos_serios_reabilitacao}. Para enfrentar esse
problema, pesquisadores têm investigado jogos que respondem ao esforço
muscular do paciente: se o esforço produz efeito visível e imediato na
tela, o exercício deixa de ser tarefa abstrata e passa a ter retorno
concreto~\cite{ref:biofeedback_engajamento}.

A eletromiografia de superfície é uma das formas mais comuns de captar
esse esforço. Ela mede, por eletrodos colocados sobre a pele, a
atividade elétrica que acompanha a contração
muscular~\cite{ref:semg_fundamentos}. Por não exigir agulhas nem
procedimentos invasivos, é adequada ao uso repetido. Aplicada aos
músculos da face, permite acompanhar quadros como a paralisia facial
periférica e sequelas de acidente vascular encefálico, nos quais a
redução da força de contração e a assimetria entre os lados do rosto
são sinais clínicos relevantes~\cite{ref:semg_facial_clinico}.

É nesse contexto que se insere o EMGInvaders, sistema desenvolvido
por~\textbf{[PREENCHER: citar o artigo do próprio projeto, se já
submetido; caso contrário, ``desenvolvido pelos autores'']}. Trata-se
de um jogo inspirado no Space Invaders, no qual a contração dos
músculos zigomáticos maiores --- responsáveis por elevar os cantos da
boca --- controla a ação em tela. O sinal é captado por dois eletrodos,
um de cada lado do rosto, com eletrodo de referência na testa.

\subsection{Problema}

Sistemas desse tipo enfrentam uma dificuldade que não é de programação,
mas da natureza do sinal: \textbf{a relação entre o esforço do paciente
e o valor medido não se mantém estável ao longo do tempo}. Basta que os
eletrodos sejam recolocados alguns milímetros adiante, que a pele
esteja mais ou menos hidratada, ou que o paciente esteja mais cansado,
para que a mesma contração produza leitura
diferente~\cite{ref:variabilidade_intersessao}.

A consequência prática é direta. O ajuste feito hoje não serve para
amanhã. A cada sessão é preciso repetir a calibração antes que o jogo
funcione. Em ambiente clínico, esse tempo compete com o tempo de
terapia, e quando o preparo consome parcela expressiva da sessão o
sistema tende a ser abandonado.

Há ainda um segundo efeito, dentro da própria sessão. À medida que o
músculo se cansa, a amplitude do sinal diminui de forma
progressiva~\cite{ref:fadiga_semg}. Se o sistema usa valor de corte
fixo, torna-se gradualmente menos sensível: o paciente precisa se
esforçar cada vez mais para obter a mesma resposta, justamente quando
já está fatigado.

\subsection{Por que Guardar a Calibração Não Basta}

Uma solução imediata seria salvar a calibração de cada paciente e
recarregá-la na sessão seguinte. Essa solução, porém, é inadequada por
uma razão que é específica do contexto de reabilitação e que orienta
todo este trabalho.

\textbf{Em reabilitação, o sinal deve mudar.} O objetivo do tratamento
é que o paciente recupere força e controle motor, e essa recuperação se
manifesta como alteração progressiva do padrão de
ativação~\cite{ref:recuperacao_motora_semg}. Um sistema que congela a
calibração da primeira sessão trabalha contra a própria finalidade
terapêutica: ele passa a medir o paciente que chegou, e não o paciente
que está se recuperando.

Surge daí uma tensão que não se resolve escolhendo um dos lados. Se o
sistema se adapta demais, perde a referência do que já observou e não
consegue mais reconhecer o participante nem registrar sua evolução. Se
não se adapta o suficiente, torna-se progressivamente inadequado. Essa
é exatamente a formulação do dilema entre plasticidade e estabilidade
estudado em aprendizado contínuo~\cite{ref:aprendizado_continuo}, e é
por essa razão que a proposta aqui apresentada adota esse referencial
em vez de tratar o problema como simples persistência de dados.

Há ainda uma dificuldade adicional, e é a mais interessante do ponto de
vista científico: \textbf{as duas fontes de mudança se confundem}.
Quando o sinal de hoje difere do de ontem, não é evidente se o eletrodo
foi recolocado em posição distinta ou se o paciente melhorou. A
primeira é ruído de medição e deve ser compensada; a segunda é o
resultado que se deseja medir e deve ser preservada. Um sistema que não
as separa ou apaga o progresso terapêutico ou o confunde com artefato.

\subsection{Lacuna}

A literatura reconhece a variabilidade do sinal de sEMG como obstáculo.
A abordagem predominante, contudo, é tratar cada sessão como evento
isolado: calibra-se do zero, usa-se, descarta-se. A informação
acumulada em sessões anteriores do mesmo paciente não é
aproveitada.~\textbf{[VERIFICAR: confirmar pela busca bibliográfica. Se
houver trabalhos que já acumulem perfil entre sessões, a lacuna deve
ser reformulada, não forçada.]}

\subsection{Objetivos}

O objetivo geral é propor um mecanismo de aprendizado contínuo de
perfis de ativação para o EMGInvaders, capaz de reconhecer
participantes já atendidos, atualizar seu perfil a cada nova sessão sem
descartar o histórico anterior, e ajustar continuamente o limiar de
ativação durante o uso.

Objetivos específicos:

\begin{enumerate}
    \item caracterizar como o sinal de sEMG facial de um mesmo
    participante varia entre sessões no contexto de uso do EMGInvaders;

    \item propor uma representação de perfil que não dependa da
    amplitude absoluta do sinal e que, por isso, resista à variação
    causada por reposicionamento de eletrodos;

    \item propor um mecanismo de atualização de perfil que incorpore
    informação nova sem apagar a anterior, permitindo acompanhar a
    evolução do participante ao longo do tratamento;

    \item propor um mecanismo de ajuste contínuo do limiar de ativação,
    sensível à queda de amplitude provocada pela fadiga;

    \item definir um protocolo de avaliação que compare a proposta com
    a versão atual do jogo, contemplando medidas objetivas de operação,
    de retenção do perfil ao longo das sessões e de percepção de uso.
\end{enumerate}

O restante do texto organiza-se como segue. A Seção~2 apresenta os
conceitos necessários. A Seção~3 discute trabalhos próximos e delimita
a lacuna. A Seção~4 descreve a metodologia. A Seção~5 apresenta o plano
de avaliação. A Seção~6 reúne contribuições esperadas, limitações e
cronograma.

%=====================================================================
\section{Revisão da Literatura}
%=====================================================================

\subsection{Eletromiografia de Superfície}

Quando o sistema nervoso comanda uma contração, cada unidade motora
recrutada gera um potencial elétrico que se propaga pelas fibras
musculares. O eletrodo de superfície capta a soma desses potenciais
atravessando os tecidos até a pele. O sinal resultante é de baixa
amplitude e ocupa faixa de frequência que se estende aproximadamente de
20 a 450~Hz, razão pela qual a amplificação e a filtragem são partes
essenciais de qualquer sistema de
aquisição~\cite{ref:semg_fundamentos}.

A aplicação à musculatura facial tem particularidades. Os músculos da
face são finos, superficiais e próximos uns dos outros, o que aumenta a
chance de um eletrodo captar atividade de músculos vizinhos ---
fenômeno conhecido como \textit{crosstalk}. O posicionamento dos
eletrodos torna-se, assim, mais crítico do que em membros, onde os
ventres musculares são maiores e mais separados.

\textbf{[PREENCHER]} Descrever a justificativa da montagem adotada no
EMGInvaders: escolha do zigomático maior bilateral, papel do eletrodo
de referência na testa e distância entre eletrodos. Este material já
existe no levantamento anterior do projeto.

\subsection{Fontes de Mudança no Sinal}

Convém distinguir três fontes de mudança, porque cada uma exige
tratamento diferente e a proposta desta pesquisa depende dessa
separação.

A primeira é o \textbf{reposicionamento entre sessões}. Recolocar os
eletrodos altera a distância até as fibras musculares; variações de
impedância da pele e de temperatura têm efeito
semelhante~\cite{ref:variabilidade_intersessao}. A evidência disponível
indica que esses fatores afetam predominantemente a \textbf{escala} do
sinal --- a amplitude sobe ou desce como um todo --- sem descaracterizar
o formato da contração, isto é, a sequência de início, sustentação e
relaxamento.~\textbf{[VERIFICAR: esta é a premissa central do método.
Precisa estar bem sustentada, pois a Seção~4.3 depende dela.]}

A segunda é a \textbf{fadiga}, que ocorre dentro da sessão. Com o
prolongamento do esforço, a amplitude decresce progressivamente e o
conteúdo de frequência se altera~\cite{ref:fadiga_semg}. É mudança
gradual e, ao contrário do reposicionamento, previsível em direção.

A terceira é a \textbf{recuperação motora}, que ocorre ao longo de
semanas. À medida que o paciente recupera controle, seu padrão de
ativação se modifica: contrações tendem a tornar-se mais consistentes,
mais bem sustentadas e mais simétricas entre os
lados~\cite{ref:recuperacao_motora_semg}. Diferentemente das duas
primeiras, esta não é uma perturbação a ser compensada, e sim o efeito
que o tratamento busca produzir.

A distinção operacional que este trabalho explora é a seguinte: o
reposicionamento altera predominantemente a \textit{escala} do sinal,
enquanto a recuperação altera o seu \textit{formato}. Se essa distinção
se sustentar empiricamente, ela oferece um critério para separar
compensação de registro clínico.

\subsection{Detecção de Ativação}

Para transformar o sinal em comando de jogo é preciso decidir, a cada
instante, se houve contração. O procedimento usual extrai o
\textbf{envelope} do sinal --- versão retificada e suavizada, que
representa a intensidade da atividade ao longo do tempo --- e compara
esse valor a um limiar.

A escolha do limiar é o ponto delicado. Um valor fixo é simples, mas
assume condições constantes. Alternativas adaptativas atualizam o
limiar a partir de estatísticas recentes do próprio sinal,
acompanhando mudanças graduais sem intervenção
manual~\cite{ref:limiar_adaptativo}.

\subsection{Matrix Profile}

O \textit{Matrix Profile} é uma estrutura para análise de séries
temporais. Dada uma série e um comprimento de janela $m$, associa a
cada trecho de tamanho $m$ a distância até o trecho mais parecido em
outra posição da série~\cite{ref:matrix_profile_base}. Dela extraem-se
dois elementos: os \textit{motifs}, padrões que se repetem,
correspondentes aos vales do perfil; e os \textit{discords}, trechos
sem equivalente em nenhum outro ponto, que sinalizam eventos atípicos.

A propriedade decisiva para este trabalho é o uso da \textbf{distância
$z$-normalizada}. Antes de comparar dois trechos, cada um é
transformado para média zero e desvio padrão unitário. A comparação
passa a considerar apenas o \textit{formato} do trecho, ignorando sua
escala e seu nível médio.

Esse detalhe conecta as Seções~2.2 e 2.4 e fundamenta a proposta. Se o
reposicionamento de eletrodos se manifesta sobretudo como mudança de
escala, e se a distância $z$-normalizada é insensível a escala, então
uma representação construída sobre ela tende a permanecer estável sob
reposicionamento --- e, ao mesmo tempo, permanece \textit{sensível} a
alterações de formato, que é como a recuperação motora se manifesta. A
mesma propriedade que confere robustez ao ruído preserva a
sensibilidade ao sinal clínico de interesse.

Quanto ao cálculo, o STAMP é a formulação
original~\cite{ref:matrix_profile_base}. O STOMP reaproveita o cálculo
da iteração anterior e permite incorporar dados novos de forma
incremental, sem recomputar
tudo~\cite{ref:matrix_profile_stomp}. Essa característica viabiliza o
uso em sistema de operação contínua, e por isso o STOMP é o algoritmo
adotado.

\subsection{Aprendizado Contínuo}

Aprendizado contínuo trata do problema de incorporar informação nova ao
longo do tempo sem descartar o que já foi aprendido. A dificuldade
central é conhecida como \textbf{esquecimento catastrófico}: ao ajustar
um modelo a dados recentes, o desempenho sobre dados antigos degrada de
forma abrupta~\cite{ref:esquecimento_catastrofico}. O campo organiza-se
em torno do compromisso entre \textit{plasticidade}, a capacidade de
incorporar o novo, e \textit{estabilidade}, a de preservar o
anterior~\cite{ref:aprendizado_continuo}.

Entre as estratégias propostas, a de interesse direto aqui é o
\textbf{replay}: manter uma memória com exemplos representativos do que
já foi observado e reapresentá-los junto aos dados novos durante a
atualização~\cite{ref:replay_memoria}. A questão prática é como
selecionar o que entra nessa memória. Estratégias de amostragem
aleatória garantem representatividade estatística, mas são cegas ao
conteúdo: podem guardar trechos de repouso e descartar contrações.

A Tabela~\ref{tab:mapeamento} apresenta a correspondência entre os
conceitos do campo e os elementos do problema tratado neste trabalho.
Essa correspondência não é meramente terminológica: cada linha
identifica um mecanismo concreto do sistema proposto, detalhado na
Seção~4.

\begin{table}[htb]
\centering
\caption{Correspondência entre conceitos de aprendizado contínuo e o
contexto do EMGInvaders.}
\label{tab:mapeamento}
\small
\begin{tabular}{@{}p{3.6cm}p{8.2cm}@{}}
\toprule
\textbf{Conceito} & \textbf{Correspondência no EMGInvaders} \\
\midrule
Tarefa / regime
& Uma sessão de um participante, com determinada montagem de eletrodos
e determinado estágio de recuperação \\
\addlinespace
Mudança de regime
& Reposicionamento dos eletrodos, entre sessões ou durante uma delas \\
\addlinespace
Plasticidade
& Acompanhar a evolução do padrão de ativação ao longo do tratamento \\
\addlinespace
Estabilidade
& Continuar reconhecendo o participante e preservar o registro das
sessões anteriores \\
\addlinespace
Esquecimento
& Perder a capacidade de reconhecer sessões antigas após atualizações
sucessivas do perfil \\
\addlinespace
Memória de \textit{replay}
& Conjunto de \textit{motifs} representativos das sessões já
realizadas \\
\addlinespace
Seleção de exemplares
& Escolha por \textit{motifs} do Matrix Profile, em vez de amostragem
aleatória \\
\bottomrule
\end{tabular}
\end{table}

O último item merece destaque, por ser o ponto em que as duas
abordagens se articulam. \textit{Motifs} são, por definição, os trechos
que mais se repetem em uma série. Selecioná-los para compor a memória
significa guardar exatamente os padrões característicos de contração
daquele participante, e não trechos escolhidos ao acaso. A estrutura do
Matrix Profile fornece, assim, o critério de seleção que o mecanismo de
replay requer.

\subsection{Jogos Sérios em Reabilitação}

\textbf{[PREENCHER]} Apresentar brevemente o campo: como jogos
controlados por biosinais têm sido usados em reabilitação, que
resultados relatam quanto a adesão, e como costumam tratar a
calibração. Este último ponto prepara a Seção~3.

%=====================================================================
\section{Trabalhos Relacionados}
%=====================================================================

\textbf{[PREENCHER: de quatro a seis parágrafos.]} Organize a discussão
nos quatro eixos abaixo e encerre cada parágrafo apontando o que aquele
trabalho \textit{não} resolve do problema aqui tratado. A seção precisa
convergir para a lacuna, e não funcionar como lista de resumos.

\begin{itemize}
    \item \textbf{Jogos de reabilitação controlados por sEMG.} Como
    esses sistemas lidam com a calibração? Repetem o procedimento a
    cada uso? Registram alguma evolução do paciente?

    \item \textbf{Adaptação à variabilidade do sinal.} Trabalhos sobre
    recalibração, adaptação de domínio ou transferência entre sessões
    em sEMG. Eles acumulam informação ou apenas se ajustam ao presente?

    \item \textbf{Aprendizado contínuo aplicado a sinais fisiológicos.}
    Existem aplicações a biosinais? Como tratam o esquecimento?

    \item \textbf{Matrix Profile aplicado a sinais fisiológicos.}
    Existe aplicação a sEMG? A resposta define se a contribuição aqui é
    de aplicação ou de método --- e a ausência também é achado que
    merece registro.
\end{itemize}

A Tabela~\ref{tab:relacionados} resume a comparação.

\begin{table}[htb]
\centering
\caption{Comparação entre os trabalhos analisados e a proposta.}
\label{tab:relacionados}
\small
\begin{tabular}{@{}lcccc@{}}
\toprule
\textbf{Trabalho} & \textbf{Domínio} & \textbf{Perfil} & \textbf{Atualiza} & \textbf{Avaliação} \\
                  &                  & \textbf{entre sessões} & \textbf{sem esquecer} & \textbf{com usuários} \\
\midrule
{[PREENCHER]} & sEMG de membro & Não & Não & Sim \\
{[PREENCHER]} & sEMG facial    & Não & Não & Não \\
{[PREENCHER]} & sEMG de membro & Sim & Não & Sim \\
{[PREENCHER]} & Séries gerais  & Sim & Sim & Não \\
\midrule
\textbf{Esta proposta} & \textbf{sEMG facial} & \textbf{Sim} & \textbf{Sim} & \textbf{Sim} \\
\bottomrule
\end{tabular}
\end{table}

% ATENCAO: a ultima linha marca "Sim" em tudo. Isso so se sustenta se
% as colunas forem criterios honestos, e nao criterios escolhidos para
% que apenas voce os satisfaca. Se algum trabalho ja atualizar perfil
% entre sessoes, mantenha o "Sim" dele e reposicione sua contribuicao.
% Comparativo enviesado e o erro mais facil de uma banca identificar.

\textbf{[PREENCHER]} Encerrar com parágrafo que enuncie a lacuna de
forma direta, retomando a tabela.

%=====================================================================
\section{Proposta de Metodologia}
%=====================================================================

\subsection{Visão Geral}

A arquitetura divide-se em três camadas, apresentadas na
Figura~\ref{fig:arquitetura}. A divisão responde a uma restrição
concreta de tempo: o jogo precisa decidir a cada quadro se houve
contração, ao passo que o reconhecimento do participante, a extração de
padrões e a consolidação do perfil ocorrem em momentos específicos, sem
prejuízo da fluidez da interação.

A separação é deliberada. O cálculo do Matrix Profile tem custo
quadrático no tamanho da série e não caberia no laço do jogo. Situado
na camada de adaptação, opera em fronteiras de segmento e no
encerramento da sessão, realimentando a camada online por meio de um
único valor: o limiar atualizado.

\begin{figure}[htb]
\centering
\begin{tikzpicture}[
    box/.style={draw, rounded corners=2pt, minimum width=2.9cm,
                minimum height=0.9cm, align=center, font=\small},
    mem/.style={draw, rounded corners=2pt, minimum width=2.9cm,
                minimum height=0.9cm, align=center, font=\small,
                fill=gray!12},
    lay/.style={font=\scriptsize\itshape, gray},
    ar/.style={-{Latex[length=2mm]}, thick},
    arb/.style={-{Latex[length=2mm]}, thick, dashed}
]

\node[box] (sensor) {Sensor sEMG\\\scriptsize filtro analógico 20--450\,Hz};
\node[box, right=1.4cm of sensor] (arduino) {Arduino\\\scriptsize amostragem e envio};

\node[box, below=1.2cm of sensor] (env) {Envelope\\\scriptsize janela deslizante};
\node[box, right=1.4cm of env] (det) {Detecção\\\scriptsize limiar $\tau_t$};
\node[box, right=1.4cm of det] (game) {Jogo\\\scriptsize resposta em tela};

\node[box, below=1.6cm of env] (mp) {Matrix Profile\\\scriptsize STOMP};
\node[mem, right=1.4cm of mp] (perfil) {Memória de perfis\\\scriptsize \textit{motifs} acumulados};
\node[box, right=1.4cm of perfil] (cons) {Consolidação\\\scriptsize atualização + \textit{replay}};

\draw[ar] (sensor) -- (arduino);
\draw[ar] (arduino.south) -- ++(0,-0.35) -| (env.north);
\draw[ar] (env) -- (det);
\draw[ar] (det) -- (game);
\draw[ar] (env.south) -- ++(0,-0.5) -| (mp.north);
\draw[ar] (mp) -- (perfil);
\draw[ar] (perfil) -- (cons);
\draw[arb] (cons.south) -- ++(0,-0.5) -| (perfil.south);
\draw[ar] (cons.north) -- ++(0,0.6) -| (det.south);

\node[lay, left=0.2cm of sensor] {\rotatebox{90}{Aquisição}};
\node[lay, left=0.2cm of env] {\rotatebox{90}{Online}};
\node[lay, left=0.2cm of mp] {\rotatebox{90}{Adaptação}};

\end{tikzpicture}
\caption{Arquitetura proposta. A camada online decide a cada quadro. A
camada de adaptação extrai padrões, consulta a memória de perfis,
atualiza-a ao final da sessão (seta tracejada) e devolve o limiar
atualizado ao detector.}
\label{fig:arquitetura}
\end{figure}

\subsection{Aquisição do Sinal}

A aquisição utiliza sensor compatível com o Muscle Sensor v3, baseado
no amplificador de instrumentação AD8226, conectado a um Arduino Uno.
Convém registrar com precisão onde ocorre cada etapa do processamento:
\textbf{a filtragem na faixa de 20 a 450~Hz e a retificação são
realizadas pelo estágio analógico do próprio sensor}, antes da
conversão digital. O sinal que chega ao computador já é, portanto, um
envelope, e não o sinal bruto.

Essa distinção delimita o escopo. Análises dependentes do conteúdo de
frequência --- como a estimativa de fadiga por deslocamento espectral
--- exigiriam sinal bruto amostrado a taxas bem mais altas e estão fora
desta proposta. O envelope permite caracterizar a intensidade e o
formato temporal das contrações, que é o suficiente aqui.

\textbf{[PREENCHER]} Informar a taxa de amostragem efetiva do sistema
atual e as restrições que a determinam, no firmware e na comunicação
serial.

\subsection{Representação do Perfil}

O perfil de um participante é o conjunto de \textit{motifs} extraídos
de suas sessões anteriores --- os padrões de contração que se repetem
de forma característica quando aquela pessoa joga. A representação
guarda, para cada \textit{motif}, a forma $z$-normalizada e, em
separado, a estatística de amplitude correspondente.

Essa separação entre forma e amplitude é intencional e é o que permite
tratar de maneira distinta os dois tipos de mudança discutidos na
Seção~2.2. Como o Matrix Profile compara trechos após normalização,
duas contrações de mesmo formato são consideradas semelhantes ainda que
uma tenha amplitude muito maior. A componente de forma tende, assim, a
permanecer estável sob reposicionamento de eletrodos, enquanto a
componente de amplitude absorve essa variação e é recalibrada a cada
sessão.

O funcionamento em uma nova sessão é o seguinte. O sistema registra os
primeiros segundos de sinal, extrai os padrões presentes e os compara
com os perfis armazenados. Se a semelhança com algum perfil for
suficiente, o participante é considerado reconhecido, sua componente de
forma é reaproveitada e apenas a escala de amplitude é reajustada. Caso
contrário, executa-se a calibração convencional e cria-se novo perfil.

Três parâmetros precisam ser definidos:

\begin{itemize}
    \item \textbf{Comprimento da janela $m$.} Deve ser especificado em
    \textit{segundos}, e o valor em amostras derivado da taxa efetiva.
    Fixar $m$ em amostras faria seu significado temporal mudar
    silenciosamente caso a taxa de aquisição fosse alterada.
    \textbf{[PREENCHER: faixa a testar, por exemplo 1, 2{,}5 e 5
    segundos.]}

    \item \textbf{Número de \textit{motifs} por perfil.} Define a
    capacidade da memória. \textbf{[PREENCHER]}

    \item \textbf{Limiar de reconhecimento.} Distância abaixo da qual um
    participante é considerado reconhecido. É o parâmetro mais sensível
    da proposta: valor permissivo demais confunde participantes
    distintos; restritivo demais elimina o benefício, pois ninguém é
    reconhecido. \textbf{[PREENCHER: como será determinado.]}
\end{itemize}

\subsection{Consolidação do Perfil entre Sessões}

Ao final de cada sessão, o perfil do participante é atualizado com os
padrões recém-observados. É nesta etapa que o dilema entre plasticidade
e estabilidade se manifesta concretamente: substituir o perfil pelos
dados da sessão mais recente equivale a esquecer as anteriores, o que
comprometeria tanto o reconhecimento futuro quanto o registro da
evolução do participante.

Adota-se, por isso, um mecanismo de consolidação por \textit{replay}. O
perfil atualizado é formado a partir da união dos \textit{motifs} da
sessão nova com os \textit{motifs} já armazenados, e não apenas dos
primeiros. Quando a capacidade da memória é atingida, a seleção do que
permanece observa dois critérios:

\begin{itemize}
    \item \textbf{Representatividade.} Preferem-se os \textit{motifs}
    com menor distância no Matrix Profile, isto é, os padrões que mais
    se repetiram --- e que, portanto, melhor caracterizam a contração
    típica daquele participante.

    \item \textbf{Cobertura temporal.} Reserva-se parte da memória para
    \textit{motifs} de sessões antigas, de modo que a atualização não
    seja dominada pelas sessões recentes. Sem essa reserva, o perfil
    convergiria gradualmente para o estado atual do participante,
    perdendo o registro de sua trajetória.
\end{itemize}

A diferença em relação à amostragem aleatória, comum em mecanismos de
replay~\cite{ref:replay_memoria}, é que a seleção aqui é informada pelo
conteúdo. Um sorteio uniforme sobre a sessão guardaria
predominantemente trechos de repouso, que ocupam a maior parte do
tempo. A seleção por \textit{motifs} guarda contrações.

\textbf{[PREENCHER]} Definir a proporção da memória reservada a sessões
antigas e o critério de descarte quando a capacidade é excedida.

\subsection{Detecção de Mudança Durante a Sessão}

O reconhecimento descrito na Seção~4.3 ocorre no início da sessão e
pressupõe que as condições permaneçam válidas até o fim. Essa
pressuposição pode falhar: um eletrodo pode se deslocar ou perder
aderência no decorrer do uso, e a calibração torna-se inválida a partir
daquele instante.

Propõe-se monitorar essa possibilidade comparando continuamente os
padrões observados com os do perfil ativo. Um afastamento sustentado
--- e não um pico isolado, que pode corresponder a artefato pontual ---
indica mudança de condição~\cite{ref:deteccao_drift}. Detectada a
mudança, o sistema reancora a escala de amplitude sem descartar a
componente de forma, uma vez que esta última não deveria ter sido
afetada por deslocamento de eletrodo.

Esse mecanismo também produz um dado útil ao registro clínico: o
momento e a frequência das reancoragens indicam problemas de fixação
dos eletrodos, informação hoje indisponível ao terapeuta.

\subsection{Limiar Adaptativo}

O limiar de ativação $\tau_t$ é atualizado ao longo da sessão por média
móvel exponencial. A cada instante, a estatística de referência do
sinal combina-se com o valor anterior, de modo que observações recentes
pesam mais que antigas:

\begin{equation}
\mu_t = \alpha\, s_t + (1-\alpha)\,\mu_{t-1},
\qquad
\tau_t = \mu_t + k\,\sigma_t,
\end{equation}

\noindent onde $s_t$ é a estatística observada no instante $t$,
$\alpha \in (0,1)$ controla a velocidade de adaptação e $k$ define a
sensibilidade da detecção. Valores maiores de $\alpha$ tornam o sistema
mais reativo; menores, mais estável.

Há aqui uma decisão que não é técnica, mas clínica. Se o limiar
acompanhar integralmente a queda de amplitude causada pela fadiga, o
jogo continuará respondendo bem --- porém a fadiga desaparecerá dos
dados, e ela é frequentemente aquilo que o terapeuta deseja observar.
Propõe-se, por isso, restringir $\tau_t$ a uma faixa em torno da
calibração inicial e \textbf{registrar a deriva observada como dado da
sessão}, em vez de absorvê-la silenciosamente. O sistema permanece
utilizável sem ocultar informação clinicamente relevante.

\textbf{[PREENCHER]} Definir a amplitude dessa faixa e discutir a
escolha com a orientação clínica do projeto.

%=====================================================================
\section{Proposta de Avaliação Experimental}
%=====================================================================

\subsection{Delineamento}

A avaliação compara duas condições: a versão atual do EMGInvaders, que
exige calibração completa a cada sessão, e a versão proposta, com
perfil consolidado e limiar adaptativo. A versão atual funciona como
\textit{baseline}.

Cada participante realiza múltiplas sessões separadas por dias. Essa
característica é indispensável em dois sentidos. Sem intervalo entre
sessões não ocorre o reposicionamento de eletrodos que produz a
variabilidade a ser tratada. E sem uma sequência de sessões não é
possível avaliar se a consolidação preserva o que foi aprendido ---
questão central da proposta.

\textbf{[PREENCHER]} Número de participantes, critérios de inclusão e
exclusão, número de sessões por participante, intervalo entre elas e
controle da ordem de apresentação das condições.

\vspace{0.5em}
\noindent\textbf{[VERIFICAR --- condição para a execução]} O protocolo
aprovado sob o Parecer CEP 8.416.967 precisa contemplar dois pontos:
sessões múltiplas por participante com intervalo de dias, e
armazenamento de dados vinculados ao participante entre sessões. O
perfil consolidado exige identificar a pessoa ao longo do tempo, o que
difere de coleta anônima pontual. Caso o protocolo não cubra esses
aspectos, é necessária emenda, e o prazo deve constar do cronograma.

\subsection{Métricas}

As métricas planejadas constam da Tabela~\ref{tab:metricas}. Elas
cobrem quatro aspectos: o custo de preparação, que é o problema de
partida; a retenção do perfil ao longo das sessões, que verifica se a
consolidação cumpre seu papel; a qualidade da detecção, para confirmar
que o ganho de tempo não custou confiabilidade; e a percepção de quem
usa, determinante para a adoção em contexto de reabilitação.

\begin{table}[htb]
\centering
\caption{Métricas planejadas.}
\label{tab:metricas}
\small
\begin{tabular}{@{}p{3.0cm}p{5.6cm}p{3.4cm}@{}}
\toprule
\textbf{Aspecto} & \textbf{Métrica} & \textbf{Resultado esperado} \\
\midrule
Custo de preparação
& Tempo entre o início da sessão e o momento em que o jogo se torna
operacional
& Menor que o do \textit{baseline} \\
\addlinespace
Reconhecimento
& Proporção de sessões em que o participante é corretamente
identificado
& {[PREENCHER]} \\
\addlinespace
Retenção
& Capacidade de reconhecer sessões antigas após sucessivas
atualizações do perfil
& Sem degradação sistemática \\
\addlinespace
Efeito da consolidação
& Comparação entre atualização com \textit{replay} e substituição
simples do perfil
& Vantagem para o \textit{replay} \\
\addlinespace
Confiabilidade
& Frequência de disparos indevidos ao longo da sessão
& Não superior ao \textit{baseline} \\
\addlinespace
Comportamento sob fadiga
& Variação da resposta entre o início e o fim da sessão
& {[PREENCHER]} \\
\addlinespace
Usabilidade
& Escala aplicada ao final de cada sessão
& {[PREENCHER: definir o instrumento]} \\
\bottomrule
\end{tabular}
\end{table}

A métrica de \textbf{retenção} merece comentário, por ser a que traduz
o problema do esquecimento para este contexto. Após o perfil de um
participante ter sido atualizado pelas sessões 2 a $n$, verifica-se se
ele ainda reconhece o registro da sessão 1. Degradação progressiva
indicaria que a consolidação não está preservando o histórico, com duas
consequências: perda de confiabilidade do reconhecimento e perda do
registro longitudinal de evolução.

\textbf{[PREENCHER]} Para cada métrica, definir como será medida na
prática. O tempo até o sistema tornar-se operacional, por exemplo,
requer critério de parada explícito. Quanto à usabilidade, verificar se
o instrumento escolhido possui versão validada em português.

\subsection{Análise dos Resultados}

Estudos dessa natureza costumam envolver poucos participantes, o que
limita o uso de testes que pressupõem distribuição normal. Prevê-se,
portanto, o uso de testes não paramétricos, acompanhados do relato do
tamanho do efeito --- mais informativo que a simples significância
quando a amostra é pequena.

O plano de análise é declarado antes da coleta, de modo a evitar que a
escolha do teste seja influenciada pelos dados observados.

\textbf{[PREENCHER]} Especificar os testes previstos para cada métrica.

%=====================================================================
\section{Considerações Finais e Cronograma}
%=====================================================================

\subsection{Contribuições Esperadas}

Espera-se que o trabalho produza quatro contribuições. A primeira é a
formulação da calibração inter-sessão em sEMG como problema de
aprendizado contínuo, com o mapeamento explícito apresentado na
Tabela~\ref{tab:mapeamento}. A segunda é um mecanismo de seleção de
exemplares baseado em \textit{motifs}, alternativa informada pelo
conteúdo à amostragem aleatória usual em replay. A terceira é a
integração desses mecanismos a um sistema em funcionamento, com
avaliação em condições de uso. A quarta é um protocolo de avaliação
replicável, que inclui métricas de retenção adaptadas ao contexto de
reabilitação.

\subsection{Limitações}

Algumas limitações são conhecidas desde o planejamento e convém
declará-las.

A amostra tende a ser reduzida e composta por participantes sem a
condição clínica alvo. Isso restringe especialmente a avaliação do
componente de plasticidade: sem participantes em recuperação efetiva,
não é possível observar a alteração de padrão que a proposta pretende
acompanhar. A avaliação, nesta etapa, limita-se à estabilidade e à
retenção, ficando a plasticidade como objeto de estudo posterior.

A aquisição restringe-se ao envelope do sinal, o que exclui análises
baseadas em conteúdo de frequência. A montagem emprega apenas dois
canais, o que reduz o detalhamento possível da atividade facial. Por
fim, não se identificou base pública de sEMG facial com múltiplas
sessões que permitisse validação externa do
método~\textbf{[VERIFICAR na busca bibliográfica]}.

\subsection{Cronograma}

A Figura~\ref{fig:cronograma} apresenta o cronograma em horizonte de
seis meses. A regularização ética aparece logo no início e em paralelo
à revisão porque constitui caminho crítico: sem ela, a coleta em
múltiplas sessões não pode ocorrer, e a avaliação fica inviabilizada.

\begin{figure}[htb]
\centering
\resizebox{\textwidth}{!}{
\begin{ganttchart}[
    hgrid, vgrid,
    x unit=1.6cm,
    y unit title=0.6cm,
    y unit chart=0.56cm,
    title label font=\small,
    bar label font=\small,
    group label font=\small\bfseries,
    milestone label font=\small\itshape,
    bar/.append style={fill=gray!45},
    group/.append style={fill=gray!75}
]{1}{6}
\gantttitle{Mês}{6} \\
\gantttitlelist{1,...,6}{1} \\

\ganttgroup{Fundamentação}{1}{2} \\
\ganttbar{Revisão da literatura}{1}{2} \\
\ganttbar{Regularização ética}{1}{2} \\

\ganttgroup{Desenvolvimento}{2}{4} \\
\ganttbar{Instrumentação da aquisição}{2}{3} \\
\ganttbar{Representação de perfil}{3}{4} \\
\ganttbar{Consolidação e limiar}{3}{4} \\
\ganttbar{Integração ao jogo}{4}{4} \\

\ganttgroup{Avaliação}{4}{6} \\
\ganttbar{Coleta em múltiplas sessões}{4}{5} \\
\ganttbar{Análise dos resultados}{5}{6} \\

\ganttgroup{Redação}{5}{6} \\
\ganttbar{Escrita do artigo}{5}{6} \\
\ganttmilestone{Submissão}{6}
\end{ganttchart}
}
\caption{Cronograma estratégico de execução.}
\label{fig:cronograma}
\end{figure}

%=====================================================================
\bibliographystyle{plain}
\bibliography{referencias}
%=====================================================================

%=====================================================================
\section*{Declaração sobre o Uso de Inteligência Artificial Generativa}
%=====================================================================

\textbf{[PREENCHER]} Descrever com precisão em que etapas houve uso de
IA generativa, qual ferramenta, e o que foi verificado pelos autores.
Declarações genéricas do tipo ``a IA foi usada para revisar o texto''
são frágeis. As referências bibliográficas, em particular, devem ser
declaradas como localizadas e conferidas manualmente pelos autores ---
e precisam de fato tê-lo sido.

\end{document}
